
\uline{par} @@M0@@. \uline{Si} @@M1@@ \uline{est le morphisme structural de} @@M2@@, \uline{alors} @@M3@@, \uline{défini car} @@M4@@ \uline{est plat}, \uline{est surjectif et définit un isomorphisme} :
@@M5@@

Soient @@M6@@, @@M7@@ la base évidente de @@M8@@, @@M9@@ le sous-Module de @@M10@@ engendré par @@M11@@. Le sous-foncteur de @@M12@@ qui à tout @@M13@@-schéma @@M14@@ associe le sous-groupe des automorphismes de @@M15@@ qui laissent invariants les sous-modules @@M16@@ pour tout @@M17@@ est représentable par un sous-schéma en groupes @@M18@@ de @@M19@@, qui est un sous-groupe de Borel de @@M20@@. Le faisceau quotient @@M21@@ (pour la topologie fpqc) est représentable par le schéma des drapeaux de @@M22@@, qui sera noté @@M23@@. Soient @@M24@@ les quotients canoniques de l'image inverse de @@M25@@ sur @@M26@@, @@M27@@ le noyau de @@M28@@, et @@M29@@.

Par ailleurs, @@M30@@ est produit semi-direct de son radical unipotent @@M31@@ et de son sous-tore maximal @@M32@@ laissant les @@M33@@ invariants. Soit @@M34@@. Alors @@M35@@ est un torseur de base @@M36@@ sous le groupe @@M37@@. Si @@M38@@ est un @@M39@@-schéma, la donnée d'un @@M40@@-point de @@M41@@ équivaut à la donnée, sur un @@M42@@ fidèlement plat quasi-compact sur @@M43@@, d'un automorphisme de @@M44@@ modulo les unipotents de @@M45@@, c'est-à-dire d'un drapeau de @@M46@@, les sous-modules du drapeau étant libres, et pour tout @@M47@@, d'une classe d'éléments @@M48@@ libres sur @@M49@@ tels que @@M50@@, soit encore à la donnée d'un drapeau de @@M51@@ et pour tout @@M52@@ d'une base de @@M53@@, cette donnée étant soumise à la condition de recollement sur @@M54@@; par descente fpqc, cela revient à la même donnée sur @@M55@@. On voit alors immédiatement que @@M56@@ est le @@M57@@-schéma affine défini par l'algèbre :
